\chapter*{Conclusion générale et perspectives}
\addcontentsline{toc}{chapter}{Conclusion générale et perspectives}
\markboth{CONCLUSION GÉNÉRALE}{}

Ce projet de fin d'année a abouti à la conception et à la réalisation de
\textbf{\montitre}, une plateforme de Business Intelligence \textbf{généraliste}
et \textbf{automatisée}. Partant du constat que la chaîne décisionnelle
traditionnelle reste manuelle, coûteuse et spécialisée, nous avons proposé une
architecture originale qui réconcilie la rigueur d'un \textbf{moteur
déterministe} et le discernement d'une \textbf{équipe de sept agents
d'intelligence artificielle}, le tout orchestré par le \textbf{protocole MCP}.

Les objectifs fixés en introduction ont été atteints. Une architecture
multi-agents à contrôle d'accès par rôle a été mise en place, et un moteur
adaptatif profile désormais n'importe quel jeu de données pour en dériver
automatiquement indicateurs, graphiques et insights. Le système produit, à
chaque exécution, un tableau de bord interactif et un rapport décisionnel
(web et PDF) reproductibles, tandis que des mécanismes de fiabilité — filets
de sécurité, vérification des chiffres et suivi des coûts — renforcent la
confiance dans les résultats. L'étude de cas menée sur 17 769 matchs
internationaux a enfin validé l'approche de bout en bout. Au-delà du produit,
ce travail nous a permis de consolider des compétences transversales :
conception logicielle en couches, intégration de modèles de langage et d'outils
via un protocole standard, ingénierie de la fiabilité face à des composants non
déterministes, et communication scientifique.

Plusieurs pistes d'évolution prolongent naturellement ce travail. L'ajout d'un
analyste conversationnel permettrait de poser des questions en langage naturel
et d'obtenir réponses et graphiques à la volée, en réutilisant le registre
d'outils MCP existant. L'ingestion pourrait être étendue à de nouvelles sources
de données — bases SQL, feuilles de calcul en ligne, stockage objet — ainsi
qu'aux jointures multi-sources. Des analyses prédictives, telles que la
prévision de séries temporelles et l'identification de facteurs explicatifs,
compléteraient utilement l'analyse descriptive actuelle. Enfin,
l'industrialisation de la plateforme — conteneurisation, persistance des
exécutions, authentification et déploiement multi-utilisateurs — ouvrirait la
voie à une mise en production.

En définitive, \textbf{\montitre} démontre qu'il est possible d'automatiser la
chaîne décisionnelle complète sans sacrifier la fiabilité, en assignant à
chaque composant — moteur déterministe et agents IA — le rôle où il excelle.
Cette complémentarité nous semble une voie prometteuse pour l'avenir des
outils décisionnels.
